% Evans 7.4: printed pp.409--420 / English PDF pp.426--437.
% Chinese: printed pp.289--296 / scan pp.307--314.
\Needspace{55mm}
\section{Semigroup Theory}
\label{sec:7-4}
\Needspace{55mm}
\subsection{Strong continuity, the generator, and its domain}
Let $X$ be a real Banach space. A strongly continuous semigroup is a family
$S(t)\in\mathcal L(X)$, $t\geq0$, satisfying
\[
 S(0)=I,\quad S(t+s)=S(t)S(s),\quad S(t)x\longrightarrow x
 \text{ in }X\text{ as }t\downarrow0.
\]
We first assume $\|S(t)\|\leq1$ (a contraction semigroup).
Its generator is the generally unbounded operator
\[
 D(G)=\left\{x:\lim_{h\downarrow0}\frac{S(h)x-x}{h}
          \text{ exists in }X\right\},\qquad
 Gx=\lim_{h\downarrow0}\frac{S(h)x-x}{h}.
\]
We use $G$ instead of the source's $A$, to distinguish it from the principal
coefficient matrix. For $u_t+Lu=0$, the generator is $G=-L$.
\Needspace{55mm}
\PDEAnnotation{强连续不等于范数连续}{要求每个固定 $x$ 的 $\|S(t)x-x\|\to0$。若要求 $\|S(t)-I\|\to0$，就排除了许多无界生成元，例如 $L^2$ 上的热半群。}

\Needspace{55mm}
\begin{theorem}[Generator properties; source \S7.4.1, Theorems 1--2]
\label{thm:7-generator}
The domain $D(G)$ is dense, and $G$ is closed. For $x\in D(G)$,
$S(t)x\in D(G)$, $GS(t)x=S(t)Gx$, and
\[
 \frac d{dt}S(t)x=GS(t)x=S(t)Gx.
\]
At $t=0$ this is a right derivative. In particular
$S(t)x-x=\int_0^tS(s)Gx\,ds$.
\end{theorem}
\begin{expandedproof}[domain invariance, density, and graph closure]
Commuting $S(t)$ with $S(h)$ and passing to the defining difference quotient
proves invariance and commutation. The right derivative follows likewise;
for the left derivative at $t>0$ use
\[
 \frac{S(t)x-S(t-h)x}{h}=S(t-h)\frac{S(h)x-x}{h}
 \longrightarrow S(t)Gx.
\]
Strong continuity and the contraction bound justify the last limit.
For arbitrary $x\in X$, let $x_t=t^{-1}\int_0^tS(s)x\,ds$.
A change of integration interval gives
\[
 x_t\in D(G),\qquad Gx_t=\frac{S(t)x-x}{t},\qquad x_t\to x.
\]
Thus the domain is dense. To prove closedness, take $x_k\to x$,
$Gx_k\to y$ in $X$ and pass to the limit in
$S(t)x_k-x_k=\int_0^tS(s)Gx_k\,ds$.
The contraction bound makes convergence of the integrals uniform on
bounded time intervals. Divide the limit identity by $t$ and let
$t\downarrow0$: its right side tends to $y$. Hence $x\in D(G)$ and $Gx=y$.
\end{expandedproof}
\Needspace{55mm}
\PDEAnnotation{闭算子与有界算子的区别}{闭性是 $x_k\to x$ 且 $Gx_k\to y$ 时才推出 $Gx=y$，不表示 $\|Gx\|\leq C\|x\|$。热方程的高频特征函数说明后一不等式不成立。}

\Needspace{55mm}
\subsection{Resolvents as Laplace transforms}
For a closed densely defined operator $G$, let $\rho(G)$ be the set of
real $\lambda$ for which $\lambda I-G:D(G)\to X$ is bijective, and
$R_\lambda=(\lambda I-G)^{-1}$. The closed graph theorem makes
$R_\lambda:X\to X$ bounded. Its range is $D(G)$.
\Needspace{55mm}
\begin{theorem}[Resolvent identities and representation; source Theorem 3]
\label{thm:7-resolvent}
For $\lambda,\mu\in\rho(G)$,
\begin{equation}\label{eq:7-resolvent-id}
 R_\lambda-R_\mu=(\mu-\lambda)R_\lambda R_\mu,\qquad
 R_\lambda R_\mu=R_\mu R_\lambda.
\end{equation}
For a contraction generator every $\lambda>0$ is in $\rho(G)$, and
\begin{equation}\label{eq:7-resolvent-laplace}
 R_\lambda x=\int_0^\infty e^{-\lambda t}S(t)x\,dt,
 \qquad\|R_\lambda\|\leq\lambda^{-1}.
\end{equation}
The integral is a norm-convergent Bochner integral.
\end{theorem}
\begin{expandedproof}[the inverse on both sides]
For the first identity insert
$(\mu I-G)-(\lambda I-G)=(\mu-\lambda)I$ between $R_\lambda$ and
$R_\mu$. These compositions are defined because each resolvent maps into
$D(G)$. Interchanging $\lambda,\mu$ gives commutation (equal parameters
are immediate).
Let $J_\lambda x$ denote the integral in \eqref{eq:7-resolvent-laplace}.
It converges since its integrand has norm at most $e^{-\lambda t}\|x\|$.
Changing variables gives
\[
 \frac{S(h)J_\lambda x-J_\lambda x}{h}
 =\frac{e^{\lambda h}-1}{h}J_\lambda x
 -\frac{e^{\lambda h}}h\int_0^he^{-\lambda t}S(t)x\,dt.
\]
The limit is $\lambda J_\lambda x-x$. Thus $J_\lambda x\in D(G)$ and
$(\lambda I-G)J_\lambda x=x$.
For $x\in D(G)$, differentiate $e^{-\lambda t}S(t)x$ and integrate:
\[
 \int_0^\infty e^{-\lambda t}S(t)(\lambda x-Gx)\,dt=x.
\]
The endpoint at infinity vanishes by the contraction bound. Therefore
$J_\lambda$ is also a left inverse, proving bijectivity and the formula.
The norm bound follows by integration. This proof never passes an
unbounded operator through an integral without a domain argument;
Problem 13 gives the equivalent graph-norm justification.
\end{expandedproof}
\Needspace{55mm}
\PDEAnnotation{预解算子的符号}{$G=-L$ 时 $R_\lambda=(\lambda I+L)^{-1}$。椭圆能量估计给出的正参数正好对应发展方程的稳定方向；若误写成 $L-\lambda I$，会引入不应出现的共振。}

\PDEMarginDiagram{7.M}{同一模态的两种衰减}
{热方程取 $E=\frac12\|u\|_2^2$，图示 $k=1$ 模态。波动方程取 $E=\frac12(\|u_t\|_2^2+\|u_x\|_2^2)$，同一模态总能量守恒。}
{ch07-energy-margin}

\Needspace{55mm}
\subsection{Constructing the semigroup from resolvent bounds}
\Needspace{55mm}
\begin{theorem}[Hille--Yosida, contraction case; source \S7.4.2, Theorem 4]
\label{thm:7-hille-yosida}
A closed densely defined linear $G$ generates a contraction semigroup if
and only if
\[
 (0,\infty)\subset\rho(G),\qquad\|R_\lambda\|\leq\lambda^{-1}
 \quad(\lambda>0).
\]
\end{theorem}
\begin{expandedproof}[Yosida approximation]
Necessity is Theorem~\ref{thm:7-resolvent}. Conversely define the bounded
operators
\[
 G_\lambda=\lambda^2R_\lambda-\lambda I=\lambda GR_\lambda.
\]
For $x\in D(G)$,
$\lambda R_\lambda x-x=R_\lambda Gx\to0$.
The bound $\|\lambda R_\lambda\|\leq1$ and density extend this convergence
to every $x\in X$. Consequently $G_\lambda x=\lambda R_\lambda Gx\to Gx$
for $x\in D(G)$.
Put $S_\lambda(t)=e^{tG_\lambda}$. Its power series gives
\[
 \|S_\lambda(t)\|\leq
 e^{-\lambda t}\sum_{k=0}^\infty
       \frac{(\lambda^2t)^k\|R_\lambda\|^k}{k!}\leq1.
\]
The resolvent identity shows that all $G_\lambda$ commute. Differentiate
$S_\mu(t-s)S_\lambda(s)x$ in $s$ and integrate to obtain
\[
 \|S_\lambda(t)x-S_\mu(t)x\|
 \leq t\|(G_\lambda-G_\mu)x\|\longrightarrow0
 \quad(x\in D(G)).
\]
Thus $S_\lambda(t)x$ converges uniformly on bounded time intervals.
Approximation of any $x\in X$ by points of $D(G)$ extends this convergence
to all $X$. The limit $S(t)$ is a contraction, is strongly continuous,
and satisfies the semigroup law: in the product limit use the uniform
operator bound to control one factor and strong convergence for the other.

Let $K$ be its generator. For $x\in D(G)$ pass to the limit in
\[
 S_\lambda(t)x-x=\int_0^tS_\lambda(s)G_\lambda x\,ds
\]
to get $S(t)x-x=\int_0^tS(s)Gx\,ds$. Dividing by $t$ proves
$G\subset K$. This inclusion alone is insufficient.
Fix $\lambda>0$ and $y\in D(K)$, and define
$x=R_\lambda(\lambda I-K)y\in D(G)$. Since $Kx=Gx$,
$(\lambda I-K)x=(\lambda I-K)y$. The operator $\lambda I-K$ is
injective by Theorem~\ref{thm:7-resolvent}, so $y=x\in D(G)$.
Thus $D(K)=D(G)$ and $K=G$.
\end{expandedproof}
\Needspace{55mm}
\PDEAnnotation{不能只证生成元的包含关系}{极限积分公式只给出 $G\subset K$。再用共同正参数下的预解算子满射和单射，才能排除 $K$ 有更大定义域，完成 $G=K$。}
For $\|S(t)\|\leq e^{\omega t}$, apply the theorem to $G-\omega I$.
The resolvent condition becomes
$\|R_\lambda\|\leq(\lambda-\omega)^{-1}$ for $\lambda>\omega$.
The theorem here has leading constant one; a general bound
$M e^{\omega t}$ requires the corresponding bounds on all resolvent powers.

\Needspace{55mm}
\subsection{Parabolic and wave generators}
Assume again smooth time-independent elliptic coefficients and smooth
bounded $U$. Set $D(G)=H^2(U)\cap V$, $Gu=-Lu$ on $H$.
\Needspace{55mm}
\begin{theorem}[Parabolic semigroup; source \S7.4.3, Theorem 5]
This $G$ generates a semigroup on $H$ with $\|S(t)\|\leq e^{\gamma t}$,
where $B[u,u]\geq\alpha\|u\|_V^2-\gamma\|u\|_H^2$.
\end{theorem}
\begin{expandedproof}[the elliptic resolvent and graph estimate]
Density follows from $C_c^\infty(U)\subset D(G)$.
If $u_k\to u$ and $Gu_k\to z$ in $H$, the elliptic estimate
\[
 \|u_k-u_\ell\|_{H^2}\leq C(\|G(u_k-u_\ell)\|_H+\|u_k-u_\ell\|_H)
\]
shows convergence in $H^2$. The trace is still zero and $Gu=z$; thus
$G$ is closed. For $\lambda>\gamma$, Lax--Milgram solves
$B[u,v]+\lambda(u,v)=(f,v)$, and elliptic regularity puts $u$ in $D(G)$.
Testing with $u$ gives
$(\lambda-\gamma)\|u\|_H^2\leq\|f\|_H\|u\|_H$.
Therefore $\|R_\lambda\|\leq(\lambda-\gamma)^{-1}$, and the shifted
Hille--Yosida theorem applies.
For $g\in D(G)$ the orbit is a classical $H$-valued solution. For $g\in H$
approximate by $D(G)$ data and use the energy estimate; the semigroup
orbit agrees with the weak parabolic solution. Domain regularity is not
claimed at $t=0$ for arbitrary $H$ data.
\end{expandedproof}
\Needspace{55mm}
\PDEAnnotation{自主与非自主发展}{单参数半群依赖系数不含时间。$L=L(t)$ 时通常要用双参数传播族 $U(t,s)$；不能把 $e^{-tL(t)}$ 当成解算子。Galerkin 法仍可直接处理此情形。}

For the wave application require $b=0$, $A=A^T$, $c\geq0$.
The symmetric form $B$ is coercive on $V$. On
$X=V\times H$ use
\[
 \|(u,v)\|_X^2=B[u,u]+\|v\|_H^2,\qquad
 \mathcal G(u,v)=(v,-Lu),
\]
with $D(\mathcal G)=(H^2(U)\cap V)\times V$.
\Needspace{55mm}
\begin{theorem}[Wave semigroup; source Theorem 6]
The operator $\mathcal G$ generates a contraction semigroup on $X$.
In fact it extends to an isometric group.
\end{theorem}
\begin{expandedproof}[wave resolvent]
The domain is dense. Closedness follows as above: convergence of
$(u_k,v_k)$ and $(v_k,-Lu_k)$ in $V\times H$ implies $v\in V$ and
$u\in H^2\cap V$, with the correct graph limit.
For $\lambda>0$, solving $(\lambda I-\mathcal G)(u,v)=(F,H_0)$ is
 equivalent to
\[
 (L+\lambda^2)u=H_0+\lambda F,\qquad v=\lambda u-F.
\]
Lax--Milgram and elliptic regularity give the unique domain solution.
Moreover $\langle\mathcal G(u,v),(u,v)\rangle_X
=B[v,u]-(Lu,v)=0$ by symmetry. Hence
\[
 \lambda\|(u,v)\|_X^2
 =\langle(F,H_0),(u,v)\rangle_X,
\]
which proves $\|R_\lambda\|\leq\lambda^{-1}$.
Hille--Yosida gives the contraction semigroup.
The same argument works for $-\mathcal G$. The two semigroups are
inverses: on domain data differentiate their product, then use density.
Both are contractions, so each preserves the norm. This proves the group
assertion as an editorial supplement, and explains energy conservation.
\end{expandedproof}
\Needspace{55mm}
\PDEAnnotation{波动的压缩实际是等距}{无阻尼、对称且时间独立时，反向演化也压缩，因此正向演化只能保持能量。不能由“压缩半群”一词推断波的能量严格衰减。}

\Needspace{55mm}
\subsection{Forcing and mild solutions: an editorial bridge}
For $f\in L^1(0,T;X)$ the variation-of-constants formula is
\[
 u(t)=S(t)g+\int_0^tS(t-s)f(s)\,ds.
\]
It defines a continuous $X$-valued mild solution; continuity follows from
strong continuity and the integrable norm bound.
For domain-valued data with continuous graph-norm forcing,
differentiation and Theorem~\ref{thm:7-generator} verify $u'=Gu+f$.
Approximation extends the formula as a mild solution for general $g,f$.
Uniqueness for classical domain-valued solutions follows by differentiating
$S(t-s)u(s)$. For parabolic forcing in $L^2(0,T;H)$ the energy estimate
and approximation identify this formula with the weak solution. A general
mild orbit is not asserted to be differentiable for arbitrary Banach-space
forcing.
